% ===== PRÉAMBULE CANONIQUE — à coller VERBATIM en tête de CHAQUE .tex =====
\documentclass[11pt,a4paper]{article}
\usepackage[utf8]{inputenc}\usepackage[T1]{fontenc}\usepackage{lmodern}
\usepackage[french]{babel}
\usepackage{amsmath,amssymb,mathtools}\usepackage{icomma}
\usepackage{siunitx}\sisetup{locale=FR}  % \qty{}{}, \si{}, \ang{}
\usepackage{xcolor}\usepackage{colortbl}
\usepackage{tikz}\usepackage{pgfplots}\pgfplotsset{compat=1.18}
\usepackage[siunitx,europeanresistors]{circuitikz} % circuits (élec.)
\usepackage[most]{tcolorbox}\usepackage{enumitem}\usepackage{array,booktabs}
\usepackage[a4paper,margin=2.2cm]{geometry}
\usetikzlibrary{arrows.meta,calc,angles,quotes,positioning,patterns,%
  backgrounds,shapes.geometric,decorations.pathmorphing,decorations.markings,intersections}
\definecolor{primary}{RGB}{222,49,99}\definecolor{primarydark}{RGB}{150,22,55}
\definecolor{defcol}{RGB}{0,114,178}\definecolor{methcol}{RGB}{52,92,148}
\definecolor{excol}{RGB}{0,135,90}\definecolor{attncol}{RGB}{200,40,55}
\tcbset{boxbase/.style={enhanced,breakable,boxrule=0.9pt,arc=2.5pt,
  left=3mm,right=3mm,top=2.3mm,bottom=2.3mm,fonttitle=\bfseries,coltitle=white}}
% --- boîtes TUTORIEL (noms EXACTS, ne pas renommer) ---
\newtcolorbox{cle}[1][]{boxbase,colback=defcol!6,colframe=defcol,
  title={\textsc{À retenir}\ifx\relax#1\relax\relax\else\textmd{ : \itshape #1}\fi}}
\newtcolorbox{methode}[1][]{boxbase,colback=methcol!7,colframe=methcol,sharp corners,
  title={\textsc{Méthode}\ifx\relax#1\relax\relax\else\textmd{ --- \itshape #1}\fi}}
\newtcolorbox{exemplebox}[1][]{boxbase,colback=excol!5,colframe=excol,
  title={\textsc{Exemple}\ifx\relax#1\relax\relax\else\textmd{ : \itshape #1}\fi}}
\newtcolorbox{piege}[1][]{boxbase,colback=attncol!6,colframe=attncol,
  title={\textsc{Piège}\ifx\relax#1\relax\relax\else\textmd{ : \itshape #1}\fi}}
\newtcolorbox{remarque}[1][]{boxbase,colback=black!4,colframe=black!45,
  title={\textsc{Remarque}\ifx\relax#1\relax\relax\else\textmd{ : \itshape #1}\fi}}
% --- boîtes EXERCICE / CORRIGÉ (noms EXACTS, auto-comptées) ---
\newtcolorbox[auto counter]{exo}[1][]{boxbase,colback=primary!4,colframe=primary,
  title={\textsc{Exercice}~\thetcbcounter\ifx\relax#1\relax\relax\else\textmd{ --- \itshape #1}\fi}}
\newtcolorbox[auto counter]{corrige}[1][]{boxbase,colback=excol!4,colframe=excol,
  title={\textsc{Corrigé}~\thetcbcounter\ifx\relax#1\relax\relax\else\textmd{ --- \itshape #1}\fi}}
\newcommand{\vv}[1]{\vec{#1}}\newcommand{\ds}{\displaystyle}
\newcommand{\R}{\mathbb{R}}\newcommand{\N}{\mathbb{N}}\newcommand{\Z}{\mathbb{Z}}
% (un \usepackage{fancyhdr}+en-tête est possible ; il est ignoré côté web)
% =========================================================================
\begin{document}
\begin{corrige}[le cadre qui pivote (moteur)]
\begin{enumerate}[label=\alph*)]
  \item Sur les côtés \textbf{horizontaux} (haut et bas) : le courant y est \textbf{parallèle} (ou
        antiparallèle) à $\vv{B}$, donc $\theta=\ang{0}$ ou \ang{180} et $F=BIL\sin\theta=0$.
        Ces deux côtés ne subissent \textbf{aucune} force.
  \item Les côtés \textbf{verticaux} (1 et 2) sont \textbf{perpendiculaires} à $\vv{B}$ : force
        \textbf{maximale} $F=BIL$. Par la règle des trois doigts, le côté 1 (courant vers le haut) est
        poussé \textbf{dans} la feuille ($\otimes$), le côté 2 (courant vers le bas) est poussé
        \textbf{hors} de la feuille ($\odot$). Ces deux forces sont \textbf{égales et opposées} : leur
        résultante est nulle (pas de translation), mais elles forment un \textbf{couple} qui fait
        \textbf{pivoter} le cadre autour de son axe vertical. C'est le principe du \textbf{moteur
        électrique}.
\end{enumerate}
\end{corrige}

\begin{corrige}[faire léviter une barre]
À l'équilibre, la force de Laplace ($\theta=\ang{90}$) compense le poids : $F=P$, soit $BIL=mg$.
On isole $I$, avec $m=\qty{0.02}{\kilogram}$ et $L=\qty{0.25}{\metre}$ :
\[ I=\frac{mg}{BL}=\frac{0{,}02\times10}{0{,}8\times0{,}25}=\frac{0{,}2}{0{,}2}=\qty{1}{\ampere}. \]
Avec \qty{1}{\ampere} (et le bon sens de courant pour que $\vv{F}$ soit vers le haut), la barre lévite.
\end{corrige}

\begin{corrige}[deux fils, courants opposés]
\begin{enumerate}[label=\alph*)]
  \item $\dfrac{F}{L}=\dfrac{\mu_0 I_1 I_2}{2\pi d}=\dfrac{2\cdot10^{-7}\times10\times15}{0{,}01}
        =\dfrac{3\cdot10^{-5}}{0{,}01}=\qty{3e-3}{\newton\per\metre}.$
  \item Les courants sont de \textbf{sens contraires} : les fils se \textbf{repoussent}.
\end{enumerate}
\end{corrige}

\begin{corrige}[longueur effectivement plongée dans le champ]
Seule la longueur \emph{dans le champ} intervient dans la force de Laplace : on prend
$\ell=\qty{0.12}{\metre}$ (et \textbf{non} \qty{0.30}{\metre}). Comme $\theta=\ang{90}$ :
\[ F=B\,I\,\ell=0{,}4\times6\times0{,}12=\qty{0.288}{\newton}. \]
Utiliser \qty{0.30}{\metre} donnerait un résultat faux : la partie hors champ ne subit aucune force.
\end{corrige}

\end{document}
